\clearpage
\section{Étape IV \textendash{} Étude commerciale du projet}

Cette étape formule le plan de marchéage du projet. Elle se compose de deux
volets : une étude documentaire du produit et de son marché, puis le plan de
marchéage proprement dit, structuré selon les quatre politiques du marketing
mix.

\subsection{Étude documentaire du produit et de son marché}

\subsubsection{Le produit}

\begin{center}
\begin{xltabular}{\textwidth}{P{4cm}L}
\toprule
\thead{Rubrique} & \thead{Contenu} \\
\midrule
\endfirsthead
\toprule
\thead{Rubrique} & \thead{Contenu} \\
\midrule
\endhead

\textbf{Description, fonctions et utilisations}
&
Paire de bouchons d'oreilles électroniques à réduction active de bruit, conçus
exclusivement pour le sommeil. Le produit ne diffuse aucun contenu : sa seule
fonction est de supprimer le bruit ambiant pendant la nuit. Il s'utilise
allongé, y compris sur le côté, et se recharge dans un boîtier dédié. Usages
visés : logement bruyant, voisinage, conjoint qui ronfle, travail en horaires
décalés, voyage. \\

\textbf{Caractéristiques techniques}
&
Réduction active par captation du bruit, calcul de l'onde en opposition de phase
et réémission. Format ultra-plat compatible avec le sommeil latéral. Embouts
interchangeables proposés en plusieurs tailles et matériaux. Autonomie visée
d'une nuit complète avec réduction active engagée. Boîtier de recharge par
contacts, rechargeable en USB-C. Absence de Bluetooth, d'application et de
connexion. Assemblage en France. \\

\textbf{Positionnement prix, service et qualité}
&
Positionnement premium aligné sur le haut du marché, à 249 euros TTC. La qualité
perçue repose sur le confort de maintien et la durée d'autonomie plutôt que sur
la richesse fonctionnelle. Le service constitue une part du positionnement :
essai de 30 jours satisfait ou remboursé, embouts de remplacement disponibles,
garantie légale de conformité de deux ans. \\
\bottomrule
\end{xltabular}
\end{center}

\subsubsection{Le marché}

\begin{center}
\begin{xltabular}{\textwidth}{P{4cm}L}
\toprule
\thead{Rubrique} & \thead{Contenu} \\
\midrule
\endfirsthead
\toprule
\thead{Rubrique} & \thead{Contenu} \\
\midrule
\endhead

\textbf{Caractéristiques de la demande}
&
La demande est portée par la montée de la préoccupation pour le sommeil : 42 \%
des Français placent le fait de bien dormir en tête des piliers de la santé, et
36 \% se déclarent gênés par le bruit au cours de la nuit, avec 1,8 réveil
nocturne en moyenne (INSV, 2026). Le segment des dispositifs d'aide au sommeil
par le son est récent et en construction : les premiers produits à réduction
active grand public sont apparus à la fin des années 2010. \\

\textbf{Type de clientèle potentielle}
&
Consommateurs individuels, adultes, achetant pour eux-mêmes ou pour offrir.
Aucune clientèle intermédiaire n'est visée : ni grossiste, ni détaillant, ni
industriel. Trois profils se distinguent : le citadin exposé au bruit extérieur,
la personne dont le conjoint ronfle, et le travailleur en horaires décalés
devant dormir en journée. \\

\textbf{Comportement des clients}
&
\textit{Habitudes d'achat} : recherche en ligne, comparaison des tests
indépendants et des avis d'utilisateurs avant décision. Achat souvent déclenché
par un épisode de privation de sommeil.
\newline
\textit{Motivations} : retrouver un sommeil continu, préserver sa santé et sa
concentration, éviter le conflit conjugal lié au ronflement.
\newline
\textit{Freins} : le prix, très supérieur à celui des bouchons en mousse ; le
doute sur l'efficacité réelle ; la crainte de l'inconfort une nuit entière ;
la méfiance envers une marque inconnue. \\

\textbf{Concurrence existante}
&
Concurrence directe : QuietOn 4 à 259 euros, positionné comme nous sur la
réduction active sans connexion ; Soundcore Sleep A30 à partir d'environ 196
euros, adossé au groupe Anker et plus complet fonctionnellement. Concurrence
indirecte : dispositifs à masquage sonore, bouchons en mousse à moins de cinq
euros, générateurs de bruit blanc, applications. L'analyse détaillée de ces
acteurs figure à l'étape III. \\
\bottomrule
\end{xltabular}
\end{center}

\begin{hypothese}
Le volume précis du marché français des dispositifs de sommeil à réduction
active n'est pas publié de manière exploitable pour un segment aussi récent.
L'estimation de la demande adressable sera affinée par l'étude de marché
conduite à l'étape V.
\end{hypothese}

\clearpage
\subsection{Le plan de marchéage}

\subsubsection{Politique de produit}

\begin{center}
\begin{xltabular}{\textwidth}{P{4cm}L}
\toprule
\thead{Élément} & \thead{Choix retenu} \\
\midrule
\endfirsthead
\toprule
\thead{Élément} & \thead{Choix retenu} \\
\midrule
\endhead

Nom, marque et logo
&
Marque \nomentreprise{}. La démarche de recherche du nom, la conception du logo
et la charte graphique sont détaillées à l'étape VI. \\

Gamme au lancement
&
Une référence unique. Ce choix évite de doubler les coûts d'outillage et de
certification sur un premier exercice, et concentre l'effort de communication
sur un seul produit. \\

Développements envisagés
&
Élargissement progressif de la gamme après validation du premier modèle, avec
des versions positionnées à des niveaux de prix différents afin de couvrir
l'entrée et le haut de gamme. Ce développement n'intervient pas avant que le
modèle initial ait démontré sa viabilité commerciale. \\

Stylique et design
&
Design sobre et discret, sans élément lumineux ni marquage voyant, conforme à un
produit destiné à la chambre. Le format plat, imposé par le sommeil latéral,
constitue l'élément distinctif de la silhouette. \\

Normes de qualité
&
Marquage CE. Conformité aux règles de sécurité et de transport applicables aux
batteries lithium. Respect du règlement européen relatif aux batteries, dont
l'exigence d'amovibilité applicable à compter du 18 février 2027. Inscription à
la filière de traitement des déchets d'équipements électriques et
électroniques. \\

Packaging
&
Emballage compact, sans plastique superflu, conçu pour l'expédition postale
directe sans suremballage. Il présente les tailles d'embouts fournies et les
conditions de l'essai de 30 jours. \\

Produits complémentaires
&
Embouts de remplacement vendus à l'unité ou par lot, dans plusieurs tailles et
matériaux. Ils constituent un revenu récurrent et prolongent la durée de vie du
produit. \\

Services associés
&
Essai de 30 jours satisfait ou remboursé. Garantie légale de conformité de deux
ans. Assistance par courriel. Reprise des produits en fin de vie. \\
\bottomrule
\end{xltabular}
\end{center}

\paragraph{L'essai de 30 jours et le traitement des retours.}
Notre différenciation repose sur le confort de maintien, or celui-ci ne peut
être évalué qu'après plusieurs nuits d'usage. L'essai de 30 jours lève donc le
principal frein à l'achat. Il a une contrepartie : la réglementation autorise un
vendeur à refuser la rétractation d'un produit descellé pour des raisons
d'hygiène, et nous y renonçons volontairement.

Le caractère interchangeable des embouts permet de limiter cette perte. Un
produit retourné est nettoyé, doté d'un embout neuf et soumis à des tests
portant sur l'ensemble de ses fonctionnalités. Il est ensuite remis en vente,
non pas comme un produit neuf, ce que la réglementation interdit, mais sous la
mention « reconditionné en France », dont l'emploi est autorisé dès lors que les
tests et interventions sont réalisés sur le territoire national (décret
n° 2022-190 du 17 février 2022).

Chaque retour analysé alimente en outre l'amélioration du produit, en
identifiant la taille ou le matériau d'embout qui a fait défaut.

\subsubsection{Politique de prix}

\paragraph{Méthode de fixation du prix.}
Deux méthodes ont été combinées. La première, dite par les coûts, part du coût
de revient unitaire auquel s'ajoute la marge nécessaire à couvrir les charges
fixes et à dégager un résultat. La seconde, dite par la concurrence, aligne le
prix sur les références du marché. Compte tenu d'un assemblage en France plus
coûteux qu'un assemblage asiatique, et d'un marché où les prix de référence sont
déjà élevés, la méthode par la concurrence est la contrainte dominante : c'est
elle qui fixe le plafond acceptable.

\begin{center}
\begin{tabularx}{\textwidth}{LP{3.2cm}}
\toprule
\thead{Élément} & \thead{Montant} \\
\midrule
Prix de vente public TTC & 249,00 \texteuro{} \\
Taux de TVA applicable & 20 \% \\
Prix de vente HT & 207,50 \texteuro{} \\
Montant de TVA collectée & 41,50 \texteuro{} \\
Coût de revient unitaire visé (hypothèse) & 90,00 \texteuro{} \\
Marge unitaire sur coût de revient (hypothèse) & 117,50 \texteuro{} \\
Taux de marge sur prix de vente HT (hypothèse) & 56,6 \% \\
\bottomrule
\end{tabularx}
\end{center}

\begin{hypothese}
Le coût de revient unitaire de 90 euros est une hypothèse de travail. Il devra
être remplacé par le résultat du chiffrage conduit à l'étape VIII, à partir de
devis réels de bureau d'études et de sous-traitant d'assemblage. La marge
affichée en dépend directement.
\end{hypothese}

\paragraph{Stratégie et positionnement prix.}
Le prix de 249 euros place le produit juste en dessous du concurrent le plus
directement comparable, tout en restant nettement au-dessus de l'offre la plus
accessible du marché. Ce positionnement est cohérent avec notre stratégie
d'écrémage : nous ne cherchons pas le volume mais une clientèle disposée à payer
pour un produit spécialisé.

Trois éléments justifient ce niveau de prix auprès du client : le confort de
maintien, l'autonomie d'une nuit complète avec réduction active engagée, et
l'assemblage en France. Aucune remise permanente n'est pratiquée, une politique
promotionnelle régulière étant incompatible avec un positionnement premium et
avec la crédibilité du prix affiché.

\subsubsection{Politique de distribution}

\paragraph{Stratégie retenue.}
La distribution repose exclusivement sur la vente en ligne depuis notre propre
site. Ce choix a été arbitré contre l'ouverture d'un canal sur une place de
marché généraliste, qui aurait apporté du trafic mais au prix d'une commission
sur chaque vente, d'une perte de la relation client et d'une mise en
concurrence immédiate avec les marques établies sur la même page.

\paragraph{Type de canal.}
Il s'agit d'un circuit \textbf{très court}, sans aucun intermédiaire entre le
producteur et le consommateur final.

\begin{center}
\begin{tabularx}{\textwidth}{LL}
\toprule
\thead{Avantages du circuit très court} & \thead{Contraintes} \\
\midrule
Marge intégralement conservée, sans commission de distributeur.
&
Aucun trafic acquis : toute visite doit être générée par nos propres moyens. \\
Maîtrise totale du discours, de la présentation du produit et du prix affiché.
&
Notoriété nulle au lancement, alors que le client ne peut pas essayer le produit
avant d'acheter. \\
Relation client directe, indispensable au traitement des retours et à la
collecte des retours d'usage.
&
Logistique d'expédition et de retour entièrement à notre charge. \\
Aucune dépendance à un distributeur susceptible de modifier ses conditions.
&
Absence de présence physique, alors qu'un produit porté dans l'oreille appelle
un besoin d'essai. \\
\bottomrule
\end{tabularx}
\end{center}

\paragraph{Logistique.}
Les commandes sont préparées et expédiées depuis nos locaux, par voie postale
suivie. Le faible encombrement du produit rend ce mode d'expédition économique.
Les retours sont réceptionnés, testés et reconditionnés en interne.

\subsubsection{Politique de communication}

\paragraph{Stratégie d'ensemble.}
Ne disposant ni de notoriété ni d'un budget publicitaire significatif, notre
communication privilégie la preuve sur la promesse. L'objectif prioritaire n'est
pas la portée mais la crédibilité : faire tester le produit par des tiers
indépendants et recueillir des avis vérifiables.

\begin{center}
\begin{xltabular}{\textwidth}{P{3.4cm}P{2.6cm}L}
\toprule
\thead{Vecteur} & \thead{Type} & \thead{Usage prévu} \\
\midrule
\endfirsthead
\toprule
\thead{Vecteur} & \thead{Type} & \thead{Usage prévu} \\
\midrule
\endhead

Site internet de la marque & Médias, internet & Support central. Présentation du
produit, explication du fonctionnement, conditions de l'essai. \\

Référencement naturel & Médias, internet & Contenus répondant aux recherches
réelles des internautes sur le bruit nocturne et le sommeil. \\

Publicité en ligne ciblée & Médias, internet & Budget limité, activé au
lancement puis ajusté selon le coût d'acquisition constaté. \\

Réseaux sociaux & Hors médias & Publications régulières, démonstrations,
réponses aux questions. Ton informatif, jamais médical. \\

Relations presse et tests produits & Hors médias & Envoi d'exemplaires à des
rédactions et testeurs spécialisés dans le son et le sommeil. Vecteur
prioritaire, car un test indépendant vaut davantage qu'une publicité. \\

Avis clients vérifiés & Hors médias & Sollicitation systématique après la
période d'essai, publiés sans filtrage sélectif. \\

Marketing direct & Hors médias & Lettre d'information à destination des
personnes inscrites, sans achat de fichier. \\

Salons et événements & Hors médias & Participation à des salons liés au sommeil
ou au fabriqué en France, permettant l'essai physique du produit. \\
\bottomrule
\end{xltabular}
\end{center}

\paragraph{Règles de communication que nous nous imposons.}
Deux limites encadrent l'ensemble de nos messages, et elles engagent notre
responsabilité juridique.

\begin{itemize}
  \item \textbf{Aucune allégation de santé.} Le discours porte sur le confort,
        l'autonomie et la simplicité. Affirmer que le produit soigne, traite ou
        améliore une pathologie du sommeil ferait basculer celui-ci dans le
        régime des dispositifs médicaux.
  \item \textbf{Aucune allégation relative aux ondes.} L'absence de Bluetooth
        est présentée comme un choix de simplicité et d'autonomie, jamais comme
        un bénéfice sanitaire. Les autorités compétentes n'établissent pas de
        risque avéré aux niveaux d'exposition courants (ANSES, 2022), et une
        promesse contraire constituerait une pratique commerciale trompeuse.
\end{itemize}

L'emploi de la mention d'origine française est par ailleurs strictement encadré.
Nous n'utiliserons jamais la formule « 100 \% français », qui figure parmi les
abus les plus sanctionnés, et la mention « fabriqué en France » ne sera employée
qu'après vérification de la règle d'origine applicable à notre nomenclature
douanière.

Le plan média détaillé, son échéancier et son budget chiffré font l'objet de
l'étape VII.
